% Fragmento público generado desde propietarios íntegros.
% Residencia: 52.7.1.
% Fuente: ../incoming/radion_monografia_20260827/manuscrito/sections/08_complejo_electromagnetismo.tex:1-58
\noindent The 3 TRIT regimes admit complex, electromagnetic, and mechanical realizations with distinct types.\par\medskip

\subsubsection{Operator quadratic of root internal of \(-1\)}

\paragraph{Exponential representation of the elliptic regime \(J^2=-I\)}

With \(J^2=-I\),
\begin{equation}
T_{\mathrm{ell}}(x,y)=e^{-xI-yJ}
=e^{-x}(\cos y\,I-\sin y\,J).
\label{eq:elliptic-exp}
\end{equation}
Under the identification generated by \(J\),
\[
T_{\mathrm{ell}}(x,y)\longleftrightarrow e^{-x-iy}.
\]
The scalar \(x\) controls dilation; \(y\) controls oriented
rotation. The form \(a+bi\) is recognized as elasticity plus rotation,
not as two arbitrarily adjoined quantities.

\paragraph{Exponential representation of the parabolic regime \(J^2=0\)}

With \(N^2=0\), the exponential is truncated:
\begin{equation}
T_{\mathrm{par}}(x,y)=e^{-x}(I-yN).
\label{eq:parabolic-exp}
\end{equation}
The regime describes the threshold at which there is neither periodic rotation nor exponential opening. In the active temporal reading it is the Euclidean present, but it retains memory of the branches that bound it.

\paragraph{Exponential Representation of the Hyperbolic Regime \(J^2=+I\)}

With \(R^2=I\) and \(P_\pm=(I\pm R)/2\),
\begin{equation}
T_{\mathrm{hyp}}(x,y)=e^{-xI-yR}
=q_+P_++q_-P_-,
\qquad q_\pm=e^{-x\mp y}.
\label{eq:hyperbolic-exp}
\end{equation}
The elementary invariants are
\begin{equation}
q_+q_-=e^{-2x},
\qquad
\frac{q_-}{q_+}=e^{2y}.
\label{eq:q-product-ratio}
\end{equation}
The product preserves the common mode; the quotient preserves the oriented separation.

Applied to the angular pair,
\[
x=\frac{\pi A}{180},\qquad y=\frac{\pi C^*}{180}.
\]
Therefore, \(A\) is the common coordinate of the publication and \(C^*\) its oriented contrast. Only after a physical reader has been applied are they recognized as electric and magnetic poles.

